\documentclass{article}

\usepackage{../handouts}

\title{CSCI 432 Handout 12: Shortest Paths in Linearly Interpolated Graphs}
\date{2 April 2026}
\author{Name: \rule{3in}{0.15mm}}

\begin{document}
\maketitle

\begin{prob}[SP in Linearly Interpolated Graphs]
    Given two weighted, directed graphs $G_0=(V,E,\omega_0)$
    and~$G_1=(V,E,\omega_1)$, define the linear interpolation at parameter $\tau \in
    [0,1]$ to be $$G_{\tau}:=(V,E,\omega_t := \tau\omega_1 + (1-\tau)\omega_0).$$
    Given a source $s \in V$ and a sink $\tau \in V$, we want to a data structure that will efficiently
    answer the query of: for $\tau \in [0,1]$, what is the shortest path from $s$ to
    $t$?
\end{prob}

Recall Dijkstra's algorithm from our last class period.
Here, we provide pseudocode for Dijkstra's SP algorithm for easy reference (this
computes just the distance; can you adjust this to find the path too?):%
\footnote{Initial pseudocode was generated from CatChat, using
the prompt ``simple psuedocode for dijkstra in latex using the algorithmic
package (not algpseudocode)''}
\begin{algorithm}[H]
    \caption{Dijkstra’s Shortest‑Path Algorithm}
    \label{alg:dijkstra}
    \begin{algorithmic}[1]   % the optional argument sets the line‑number spacing
        \REQUIRE A weighted, directed graph $G=(V,E, \omega)$ with non‑negative
        edge weights;
        a source vertex $s\in V$
        \ENSURE Shortest‑path distances $\dist[v]$ for all $v\in V$
        % and predecessor map $prev[v]$

        \STATE Initialize $\dist$, a data structure indexed by $V$, with value $\infty$ for all
            vertices except $\dist(s)=0$
        %\STATE Initialize \prev, indexed by $V$, values null everywhere
        \STATE $Q$ min-priority queue storing the elements off $V$, with $\dist[v]$ the
        priority of $v$ for all $v$)
        \WHILE{$Q \neq \emptyset$}
            \STATE $u \leftarrow \pop(Q)$
            \FOR{each $(u,v)$ out of $u$}
                \STATE $alt \leftarrow \dist(u) + \omega(u,v)$
                \IF{$alt < \dist(v)$}
                \STATE $\dist[v] \leftarrow alt$
                %\STATE $prev[v] \leftarrow u$
                \STATE $\textsc{Decrease-Key}(Q,v,\dist[v])$
                \ENDIF
            \ENDFOR
        \ENDWHILE
        \RETURN $\dist$
    \end{algorithmic}
\end{algorithm}

\pagebreak
\section*{Understanding the Problem}

\begin{figure}
    \centering
    \begin{subfigure}[b]{0.4\textwidth}
        \includegraphics[width=.9\textwidth]{graph-0}
        \caption{$G_0$}
        \label{fig:graph-0}
    \end{subfigure}
    \qquad
    \begin{subfigure}[b]{0.4\textwidth}
        \includegraphics[width=.9\textwidth]{graph-1}
        \caption{$G_1$}
        \label{fig:graph-1}
    \end{subfigure}
    \caption{The starting and ending graphs that we linearly interpolate
    between. Here, let all edges be directed from left to right.}\label{fig:graph}
\end{figure}


\begin{enumerate}
    \item For $G_0$ and $G_1$, use Dijkstra's algorithm to compute
        a shortest path from $s$ to $t$ in each graph (denote them by $p_0$ and
        $p_1$, respectively). What are their lengths
        (denote them by $\ell_0$ and $\ell_1$, respectively).

        \textbf{Answer}
        \vspace{1in}

    \item Compute $G_{0.5}$ and find a shortest path from $s$ to $t$.  What is
        its length? (Denote the path and lengthy by $p_{0.5}$ and
        $\ell_{0.5}$).

        \textbf{Answer}
        \includegraphics[height=2in]{graph-noweights}

    \pagebreak
    \item Compute $G_{0.25}$ and find the length of a shortest path from $s$ to
        $t$.  What is its length?

        \textbf{Answer}
        \includegraphics[height=2in]{graph-noweights}

    \item For which value of $t$ are the lengths of $p_0$ and $p_1$ the same
        (Note: there is one value $t$ for which this holds; we denote this by
        $t_{same}$)?

        \textbf{Answer}
        \vspace{1in}

    \item What are the values of all paths between $s$ and $t$ at $t_{same}$. Is
        there any time that the shortest path is not $p_0$ or $p_1$? (Hint: look
        at the graphs of path length versus parameter $\tau$ for all paths from $s$
        to $t$).

        \textbf{Answer}
        \vspace{1in}

    \item Provide an example of two graphs for which there exists a $t$ for which
        $p_t \notin \{p_0,p_1\}$.

        \textbf{Answer}

\end{enumerate}


\pagebreak
\section*{The Algorithm}

\begin{figure}
    \centering
    \includegraphics{path-lengths}
    \caption{Path lengths.}
\end{figure}

\begin{enumerate}
    \item What are some observations we have about the plot above? When are the
        shortest paths not unique?

        \textbf{Answer}
        \vspace{2in}

    \pagebreak
\item Consider the linearly interpolated graph in \figref{othergraph}, along with the graph
        of the lengths of all paths from~$s$ to $t$.
        What are the shortest path representatives for each parameter $\tau$?

    \begin{figure}
        \centering
        \begin{subfigure}[b]{0.3\textwidth}
            \includegraphics[width=.9\textwidth]{othergraph-0}
            \caption{$H_0$}
            \label{fig:graph-0}
        \end{subfigure}
        ~
        \begin{subfigure}[b]{0.3\textwidth}
            \includegraphics[width=.9\textwidth]{othergraph-1}
            \caption{$H_1$}
            \label{fig:graph-1}
        \end{subfigure}
        ~
        \begin{subfigure}[b]{0.33\textwidth}
            \includegraphics[width=.9\textwidth]{otherpath-lengths}
            \caption{Path Lengths}
            \label{fig:graph-1}
        \end{subfigure}
        \caption{The starting and ending graphs that we linearly interpolate
        between. Here, let all edges be directed from left to
        right.}\label{fig:othergraph}
    \end{figure}

        \textbf{Answer}

\end{enumerate}

\end{document}
